\section{Retention deletion and provenance boundaries}\label{sec:C13:retention}
Durability and retention answer different questions. Durability concerns surviving failures under a storage contract; retention concerns how long information should remain available and to whom. An agent can be durable while enforcing short retention, or retain data indefinitely while providing weak recovery guarantees. The design should state both instead of using ``persistent'' as a substitute for either.

A retention model should distinguish raw event payloads, derived summaries, embeddings, caches, effect receipts, and learned parameters. Deleting a source document does not necessarily delete the summary or embedding derived from it. Conversely, invalidating a derived record can prevent future use without physically erasing all copies. A truthful deletion claim must identify its scope and verification boundary. Cross-service backups and model-training artifacts require different mechanisms from a local database deletion.

Let $D(x)$ denote the set of derived objects depending on source object $x$. A deletion or access-revocation event can traverse this dependency relation to invalidate downstream objects. The graph need not imply that each derivative can be perfectly unlearned from a model; it identifies objects for which the application has provenance. If provenance was not recorded, the system may not know which summaries contain the deleted fact. This is a concrete operational cost of treating generated summaries as untraceable text.

Access checks should occur when evidence is retrieved and again when it is exposed to a model or another agent. A cached result may have been authorized when created but unauthorized when reused. Partitioning by tenant and project reduces accidental mixing, but semantic records can still reveal information across boundaries if extraction merged facts from multiple principals. The memory writer must preserve ownership or refuse a merge that destroys the ability to enforce access.

Past-decision audit creates a difficult case. An organization may need to explain why an agent acted while also respecting deletion of the underlying sensitive content. A record can preserve that a decision used an object with a given identifier and policy state without retaining the full payload. Such an audit supports a narrower claim than exact reconstruction. The distinction should be reflected in the recovery interface: a deleted evidence object should produce an explicit unavailable state, not a fabricated reconstruction.

Session summaries should therefore carry retention and provenance metadata alongside their semantic content. Consolidation into cross-session memory is a new use of information, not an automatic entitlement granted by the existence of a transcript. The application must decide whether the information belongs to the user, the project, a shared team, or a short-lived task. Chapter~\ref{ch:C14} analyzes how those choices interact with semantic and episodic memory formation.

